\section{Marco teórico}

\subsection{Estacionariedad}

Un proceso estocástico $\{y_t\}$ es \textbf{estacionario en sentido débil} cuando cumple tres
condiciones: su esperanza es constante, $E[y_t] = \mu$ para todo $t$; su varianza es finita y
constante, $\operatorname{Var}(y_t) = \sigma^2$; y su autocovarianza entre dos instantes
depende únicamente de la distancia que los separa,
\begin{equation}
\gamma_k = \operatorname{Cov}(y_t, y_{t-k}),
\end{equation}
sin depender de $t$ \parencite{box2015}.

La exigencia no es formal. Los modelos ARMA estiman un número fijo de parámetros a partir de
una única realización del proceso, y esa estimación sólo tiene sentido si las propiedades del
proceso permanecen invariantes a lo largo de la muestra. Cuando la media se desplaza, el
promedio muestral no estima ningún parámetro poblacional, y las regresiones entre series no
estacionarias pueden arrojar asociaciones elevadas carentes de significado, fenómeno conocido
como regresión espuria \parencite{box2015}.

La \textbf{función de autocorrelación} normaliza la autocovarianza,
\begin{equation}
\rho_k = \frac{\gamma_k}{\gamma_0},
\end{equation}
y la \textbf{función de autocorrelación parcial} mide la correlación entre $y_t$ e $y_{t-k}$
una vez descontado el efecto de los rezagos intermedios. El comportamiento conjunto de ambas es la herramienta primaria de identificación.

\subsection{Procesos autorregresivos y de medias móviles}

Un proceso \textbf{autorregresivo de orden $p$}, denotado AR$(p)$, expresa el valor corriente
como combinación lineal de sus propios rezagos más una perturbación:
\begin{equation}
y_t = c + \phi_1 y_{t-1} + \phi_2 y_{t-2} + \cdots + \phi_p y_{t-p} + \varepsilon_t,
\end{equation}
donde $\varepsilon_t$ es ruido blanco de media nula y varianza $\sigma^2$. Su autocorrelación decae de forma exponencial o sinusoidal, mientras que la parcial se anula a
partir del rezago $p$.

Un proceso de \textbf{medias móviles de orden $q$}, MA$(q)$, expresa el valor corriente como
combinación de perturbaciones presentes y pasadas:
\begin{equation}
y_t = \mu + \varepsilon_t + \theta_1 \varepsilon_{t-1} + \cdots + \theta_q \varepsilon_{t-q}.
\end{equation}
El comportamiento es el inverso: la autocorrelación se anula a partir del rezago $q$ y la parcial decae con lentitud.

El proceso \textbf{ARMA$(p,q)$} combina ambos. Con el operador de rezago $B$, definido por
$B y_t = y_{t-1}$, se escribe de forma compacta:
\begin{equation}
\phi(B)\, y_t = \theta(B)\, \varepsilon_t,
\end{equation}
donde $\phi(B) = 1 - \phi_1 B - \cdots - \phi_p B^p$ y
$\theta(B) = 1 + \theta_1 B + \cdots + \theta_q B^q$.

\subsection{Modelos ARIMA y SARIMA}

Cuando la serie no es estacionaria en niveles pero sí lo es tras diferenciar, el modelo
incorpora un operador de diferencia $\nabla = 1 - B$. El modelo \textbf{ARIMA$(p,d,q)$} se escribe
\begin{equation}
\phi(B)\,(1-B)^d\, y_t = \theta(B)\, \varepsilon_t,
\end{equation}
donde $d$ es el número de diferencias regulares necesarias. Corresponde emplearlo cuando la
serie presenta tendencia o deriva pero carece de ciclo repetitivo.

Cuando además existe un patrón que se repite cada $s$ períodos, el modelo
\textbf{SARIMA$(p,d,q)\times(P,D,Q)_s$} agrega polinomios estacionales y una diferencia
estacional $\nabla_s = 1 - B^s$:
\begin{equation}
\Phi(B^s)\,\phi(B)\,(1-B)^d\,(1-B^s)^D\, y_t = \Theta(B^s)\,\theta(B)\, \varepsilon_t .
\end{equation}
El período $s$ se fija según la periodicidad natural del fenómeno: veinticuatro para datos
horarios con ciclo diario, siete para datos diarios con ciclo semanal
\parencite{hyndman2021}.

\subsection{Modelos de vectores autorregresivos}

Cuando varias series se determinan entre sí, el modelo \textbf{VAR$(p)$} trata a
todas como endógenas y regresa cada una sobre los rezagos de todas \parencite{sims1980}. Para
un vector $\mathbf{y}_t$ de dimensión $K$:
\begin{equation}
\mathbf{y}_t = \mathbf{c} + \mathbf{A}_1 \mathbf{y}_{t-1} + \cdots + \mathbf{A}_p \mathbf{y}_{t-p} + \boldsymbol{\varepsilon}_t,
\end{equation}
donde cada $\mathbf{A}_i$ es una matriz $K \times K$. El número total de parámetros asciende a $K(Kp+1)$, y crece con el cuadrado de $K$ y de forma
lineal en $p$; de ahí que la parsimonia pese más que en el caso univariado \parencite{sims1980}.

El modelo exige que todas las series estén definidas sobre el mismo calendario y sean
estacionarias. La \textbf{función impulso-respuesta} describe la trayectoria de cada variable
ante una perturbación unitaria en otra, y permite leer la dinámica del sistema más allá de los
coeficientes individuales.

La \textbf{causalidad de Granger} contrasta si los rezagos de una variable mejoran la
predicción de otra, más allá de lo que aportan los rezagos propios de esta última
\parencite{granger1969}. La hipótesis nula sostiene que los coeficientes asociados a los
rezagos de la variable candidata son conjuntamente nulos. El rechazo indica precedencia temporal con contenido predictivo, no causalidad en sentido
estructural.

\subsection{Contrastes de hipótesis empleados}

La Tabla~\ref{tab:contrastes} resume los contrastes aplicados a lo largo del trabajo, con sus
hipótesis y la regla de decisión correspondiente.

\begin{table}[htbp]
\centering
\caption{Contrastes empleados, con sus hipótesis y regla de decisión}
\label{tab:contrastes}
\small
\begin{tabular}{p{3.1cm}p{4.2cm}p{4.2cm}p{3.1cm}}
\toprule
\textbf{Contraste} & \textbf{Hipótesis nula} & \textbf{Hipótesis alternativa} & \textbf{Decisión al 5\%} \\
\midrule
Dickey-Fuller aumentada & Existe una raíz unitaria; la serie no es estacionaria & La serie es estacionaria & Se rechaza si $p < 0{,}05$ \\
\addlinespace
KPSS & La serie es estacionaria en torno a un nivel o tendencia & Existe una raíz unitaria & Se rechaza si $p < 0{,}05$ \\
\addlinespace
Ljung-Box & Las autocorrelaciones hasta el rezago $m$ son conjuntamente nulas & Existe autocorrelación & Se rechaza si $p < 0{,}05$ \\
\addlinespace
Jarque-Bera & Los residuos siguen una distribución normal & No siguen una distribución normal & Se rechaza si $p < 0{,}05$ \\
\addlinespace
Razón de verosimilitud & Los parámetros restringidos son conjuntamente nulos & Al menos uno es no nulo & Se rechaza si $p < 0{,}05$ \\
\addlinespace
Causalidad de Granger & Los rezagos de la variable candidata son conjuntamente nulos & Al menos uno es no nulo & Se rechaza si $p < 0{,}05$ \\
\bottomrule
\end{tabular}
\end{table}

Las pruebas de Dickey-Fuller aumentada \parencite{dickey1979} y KPSS
\parencite{kwiatkowski1992} plantean hipótesis nulas recíprocas. Esa reciprocidad justifica su
aplicación conjunta: cuando ambas coinciden, la conclusión es sólida; cuando discrepan, la evidencia muestral no basta para pronunciarse con firmeza y la decisión debe apoyarse
en elementos adicionales.

El estadístico de Ljung-Box \parencite{ljung1978} se define como
\begin{equation}
Q = n(n+2) \sum_{k=1}^{m} \frac{\hat{\rho}_k^{\,2}}{n-k},
\end{equation}
y se distribuye asintóticamente como una chi-cuadrado. Aplicado a los residuos de un modelo ajustado, es el contraste central del diagnóstico: si los residuos conservan
autocorrelación, el modelo dejó estructura sin explicar.

El contraste de razón de verosimilitud compara un modelo ajustado contra una versión
restringida mediante
\begin{equation}
LR = 2\,\bigl(\ell_{\text{completo}} - \ell_{\text{restringido}}\bigr),
\end{equation}
que se distribuye como una chi-cuadrado con tantos grados de libertad como restricciones
impuestas. De ahí se obtiene la significatividad conjunta de los parámetros.

\subsection{Criterios de selección y métricas de error}

Los criterios de información penalizan la verosimilitud según el número de parámetros $k$
sobre $n$ observaciones:
\begin{align}
\text{AIC} &= -2\ell + 2k, \\
\text{BIC} &= -2\ell + k \ln n, \\
\text{HQIC} &= -2\ell + 2k \ln(\ln n).
\end{align}
El criterio de Akaike \parencite{akaike1974} penaliza con menor severidad y tiende a
seleccionar modelos más amplios; el bayesiano \parencite{schwarz1978} impone una penalización
creciente con el tamaño muestral y favorece la parsimonia; el de Hannan-Quinn
\parencite{hannan1979} ocupa una posición intermedia. Si discrepan entre sí, el desempeño fuera de muestra rompe el empate.

Las métricas de error empleadas son el error absoluto medio, la raíz del error cuadrático
medio y el error porcentual absoluto medio:
\begin{equation}
\text{MAE} = \frac{1}{h}\sum |e_t|, \qquad
\text{RMSE} = \sqrt{\frac{1}{h}\sum e_t^2}, \qquad
\text{MAPE} = \frac{100}{h}\sum \left|\frac{e_t}{y_t}\right| .
\end{equation}
La última merece una advertencia. Al dividir por el valor observado, el error porcentual se
vuelve inestable cuando el denominador se aproxima a cero, y pierde sentido sobre series que oscilan alrededor del origen \parencite{hyndman2006}. Su empleo exige, por lo tanto, que la
serie se encuentre expresada en una escala estable y alejada de cero.

\subsection{Herramientas de cómputo}

La estimación se realizó con la biblioteca \textit{statsmodels} del lenguaje Python
\parencite{seabold2010}, que implementa los modelos SARIMAX y VAR mediante representación en
espacio de estados y estimación por máxima verosimilitud.
